\documentstyle[12pt,fullpage,steve]{article}
\setcounter{page}{1}

\begin{document}

\title{EE 150 Final Exams : Part I \\ {\bf OPEN BOOK}}
\author{Monday, December 9 1996}
\date{}
\maketitle

\begin{itemize}
\item {\bf Name:}
\item {\bf Student Id Number:}
\end{itemize} 

{\bf Read First:}
\begin{enumerate}
\item [A.] {\bf **READ** THE DIRECTIONS CAREFULLY!}
\item [B.] Any materials including your text and notes may be consulted
during the exam at your discretion. 
\item [C.] You have a maximum of 50 minutes to complete the exam.  The exam
paper will be picked up from your desk (do not get up) at the end of
the alloted time.
\item [D.] There are 2 problems worth a total of 50 points, the first
is worth 35 points and the second is worth 15.  Be careful not to
spend too much time on any particular problem.  {\bf If you get stuck,
go on to the next problem}.  You can always come back to the one that
is giving you problems.   
\item [E.] Write all answers, including any scratch paper or rough 
notes on the pages given you.
\item [F.] If you split the exam paper, write your student id number
on each and every page.
\end{enumerate}

% ---------------------------------------------------------------------------
\pagebreak
{\bf PROBLEM \#1} (35 points)

As part of your new job at a software company, your are told that you
must write a {\em function} in C which can take the address of an
array of integers (the last element of which is negative 9999), and
which returns a double equal to the average of the numbers in the
array.  

Sounds easy so far, but then you are told that you must also 
return the {\bf addresses} of the minimum value in the array, the maximum
value in the array, and the value in the array nearest to the average
value (i.e. the median).  Unfortunately your boss hates the symbols
``[`` and ``]'' and has told you to never use them in functions {\em you}
write.

{\bf Hint 1:}  If you want to define a pointer to a pointer to an
integer, you would say int **a.  This is done in the main program
below for the prototype for {\em processArray}.  You would only need
this notation in your function header.
 
{\bf Hint 2:} A {\em correct} test program and output are given below:
\begin{verbatim}
#include <stdio.h>
#define ASIZE 13

main()
{
  int  A[ASIZE] = {2, 7, 1, 3, 4, 9, 5, 6, 8, 10, 55, -19, -9999};
  int *minimum;    /* minimum points to minimum value of array A */ 
  int *maximum;    /* maximum points to maximum value of array A */ 
  int *median;     /* median  points to median  value of array A */ 
  double average;  /* average contains the average value of array A */ 

  double processArray(int *, int **, int **, int **);   

  average = processArray(A, &minimum, &maximum, &median);

  printf("The value of A[8]  is   8, and the address is %p\n", &A[8]);
  printf("The value of A[10] is -19, and the address is %p\n", &A[10]);
  printf("The value of A[11] is  55, and the address is %p\n\n", &A[11]);

  printf("Array Average = %5.2f\n", average);
  printf("Array Minimum = %i at addr %p\n", *minimum, minimum);
  printf("Array Maximum = %i  at addr %p\n", *maximum, maximum);
  printf("Array Median  = %i   at addr %p\n", *median,  median);
  return 0;
}
\end{verbatim}

The output from the given test program is shown on the next page.

% ---------------------------------------------------------------------------
\pagebreak
{\bf PROBLEM \#1} (Continued ....)

The output from the given test program is shown below.
\begin{verbatim}
The value of A[8]  is   8, and the address is f7fff530
The value of A[10] is -19, and the address is f7fff538
The value of A[11] is  55, and the address is f7fff53c
 
Array Average =  7.83
Array Minimum = -19 at addr f7fff53c
Array Maximum = 55  at addr f7fff538
Array Median  = 8   at addr f7fff530
\end{verbatim}

Write your solution below.  Explicitly state any assumptions you make.

% ---------------------------------------------------------------------------
\pagebreak
{\bf PROBLEM \#1} (This Page for Continued Solution)


% ---------------------------------------------------------------------------
\pagebreak
{\bf PROBLEM \#1} (This Page for Continued Solution)
 

% ---------------------------------------------------------------------------
\pagebreak
{\bf PROBLEM \#2} (15 points)

This new job really is a pain.  You have replaced a previous programmer
who never wrote a comment in 10 years of C programming.   You are
given a program by your boss (who is still mad about your previous
assignment) and told to {\bf re-write} it so that it is more readable, shows
some better style, has meaningful comments and does the same thing as
it does now.  He doesn't know what it does and neither does anyone
else, but it is important that you don't change the way it works!

The program you are given is as follows.
\begin{verbatim}
#include<stdio.h>
main()
{
  int *L;  int *p;
  int s, i,j, ror, rand; 

  scanf("%i",&s);
  if((L = (int *) malloc(s * s * sizeof(int))) == NULL) return -1;
  p = L;
  for(i=0;i<s;i++) 
    for(j=0;j<s;j++) {
      scanf("%i",p);
      if ( (*p < 0) || (*p > 1) ) return -2;
      p++; }
  
  p = L;
  for(i=0;i<s;i++) {
      for(j=0;j<s;j++) 
        printf("%i",*p++);
      printf("\n"); }

  p = L;
  for(i=0;i<s;i++) {
    ror  = 0;  rand = 1;
    for(j=0;j<s;j++) {
      ror  = ror  || *p;
      rand = rand && *p++; }
    printf("Row %i:  ROR = %i, RAND = %i\n", i, ror, rand); }
}
\end{verbatim}

{\bf continued on next page ...}

% ---------------------------------------------------------------------------
\pagebreak
{\bf PROBLEM \#2} (Continued ....)

For Problem 2, 
\begin{enumerate}
\item First, show the output of the program given the input:
\begin{verbatim}
3 0 1 1 1 0 1 1 1 1 
\end{verbatim}
\item And Second: re-write the program on this and the following page
including better style, comments and whatever else you think would
make the program easier to understand.  Include a comment at the beginning of
the program(s) which briefly summarizes what the program does.

\end{enumerate}

% ---------------------------------------------------------------------------
\pagebreak
{\bf PROBLEM \#2} (Continued ....)

\end{document}

/* Local Variables: */
/* compile-command: "latex ee150-final1" */
/* End: */